\chapter{Chapter06 — Exact Sparse Spectrum Learning}
%%% generated by the blueprint pipeline from FABL.Chapter06.LearningAndTesting.ExactSparseSpectrumLearning
\section{Overview}

Declarations of the Lean module \texttt{FABL.Chapter06.LearningAndTesting.ExactSparseSpectrumLearning}.

\section{Declarations}

\begin{definition}[Exact sparse spectrum sparsity bound]
\lean{FABL.exactSparseSpectrumSparsityBound}
\statementsource{Analysis-of-Boolean-Functions-by-Ryan-ODonnell}{pdf-ffbe9c7b3306-p072-b002}
\label{FABL.exactSparseSpectrumSparsityBound}
\leanok
For a natural number $k$, the exact sparse spectrum sparsity bound is $2^{k+1}$.
\end{definition}

\begin{theorem}[Lower bound on the sparsity bound]
\lean{FABL.exactSparseSpectrumSparsityBound_two_le}
\label{FABL.exactSparseSpectrumSparsityBound_two_le}
\leanok
\uses{FABL.exactSparseSpectrumSparsityBound}
\difficulty{3}
For every natural number $k$, the quantity $2^{k+1}$ is at least $2$.
\end{theorem}

\begin{definition}[Exact sparse-spectrum accuracy input]
\lean{FABL.exactSparseSpectrumAccuracyInput}
\label{FABL.exactSparseSpectrumAccuracyInput}
\leanok
\uses{FABL.SmallBiasInput}
For every positive dimension $n$ and natural number $k$, the defined accuracy input has
output dimension $n$, bias numerator $1$, and bias denominator $4\cdot 2^k$.
\end{definition}

\begin{definition}[Exact sparse spectrum Goldreich–Levin input]
\lean{FABL.exactSparseSpectrumGoldreichLevinInput}
\statementsource{Analysis-of-Boolean-Functions-by-Ryan-ODonnell}{
  pdf-ffbe9c7b3306-p083-b001,
  pdf-ffbe9c7b3306-p083-b002
}
\label{FABL.exactSparseSpectrumGoldreichLevinInput}
\leanok
\uses{FABL.DeterministicGoldreichLevinInput, FABL.exactSparseSpectrumAccuracyInput, FABL.exactSparseSpectrumSparsityBound, FABL.exactSparseSpectrumSparsityBound_two_le}
Given natural numbers $n$ and $k$ and a proof that $n>0$, take the deterministic
Goldreich–Levin input whose accuracy parameters are those prescribed by the exact sparse
spectrum construction for $n$ and $k$, and whose integral Fourier $1$-norm bound is the
exact sparse spectrum sparsity bound at $k$.
\end{definition}

\begin{theorem}[Learning parameter identity]
\lean{FABL.exactSparseSpectrumGoldreichLevinInput_learningParameter}
\label{FABL.exactSparseSpectrumGoldreichLevinInput_learningParameter}
\leanok
\uses{FABL.DeterministicGoldreichLevinInput.learningParameter, FABL.DeterministicGoldreichLevinInput.learningParameter_cast, FABL.SmallBiasInput.epsilon, FABL.degreeFourierCoefficientAccuracy, FABL.degreeFourierCoefficientAccuracy_value, FABL.exactSparseSpectrumGoldreichLevinInput}
\difficulty{4}
Let $n$ and $k$ be natural numbers with $n>0$. The learning parameter of the exact
sparse-spectrum Goldreich–Levin input equals the degree-$k+1$ Fourier coefficient
accuracy parameter.
\end{theorem}

\begin{theorem}[Granularity at the exact sparsity bound]
\lean{FABL.spectralSparsityGranularity_exactSparseSpectrumSparsityBound}
\label{FABL.spectralSparsityGranularity_exactSparseSpectrumSparsityBound}
\leanok
\uses{FABL.degreeFourierGranularity, FABL.exactSparseSpectrumSparsityBound, FABL.spectralSparsityGranularity}
\difficulty{3}
For every nonnegative integer $k$, the spectral sparsity granularity evaluated at
$2^{k+1}$ equals the degree Fourier granularity at $k+1$.
\end{theorem}

\begin{theorem}[Threshold and Fourier granularity]
\lean{FABL.exactSparseSpectrumGoldreichLevinInput_threshold_le_granularity}
\label{FABL.exactSparseSpectrumGoldreichLevinInput_threshold_le_granularity}
\leanok
\uses{FABL.DeterministicGoldreichLevinInput.familySizeBound, FABL.DeterministicGoldreichLevinInput.familySizeBound_pos, FABL.DeterministicGoldreichLevinInput.learningParameter, FABL.DeterministicGoldreichLevinInput.threshold, FABL.DeterministicGoldreichLevinInput.threshold_cast, FABL.degreeFourierCoefficientAccuracy, FABL.degreeFourierGranularity, FABL.degreeFourierGranularity_pos, FABL.exactSparseSpectrumGoldreichLevinInput, FABL.exactSparseSpectrumGoldreichLevinInput_learningParameter}
\difficulty{4}
For natural numbers $n$ and $k$ with $n>0$, form the deterministic Goldreich–Levin input
with accuracy $1/(4\cdot 2^k)$ and Fourier bound $2^{k+1}$. Its threshold is at most the
degree Fourier granularity $2(1/2)^{k+1}$.
\end{theorem}

\begin{theorem}[Accuracy below half the Fourier granularity]
\lean{FABL.exactSparseSpectrumGoldreichLevinInput_coefficientEpsilon_lt_half_granularity}
\label{FABL.exactSparseSpectrumGoldreichLevinInput_coefficientEpsilon_lt_half_granularity}
\leanok
\uses{FABL.DeterministicGoldreichLevinInput.coefficientInput, FABL.DeterministicGoldreichLevinInput.coefficientInput_epsilon, FABL.DeterministicGoldreichLevinInput.learningParameter, FABL.DeterministicGoldreichLevinInput.threshold, FABL.SmallBiasFourierInput.epsilon, FABL.degreeFourierCoefficientAccuracy, FABL.degreeFourierGranularity, FABL.degreeFourierGranularity_pos, FABL.exactSparseSpectrumGoldreichLevinInput, FABL.exactSparseSpectrumGoldreichLevinInput_learningParameter, FABL.goldreichLevinActiveCap, FABL.goldreichLevinActiveCap_pos}
\difficulty{4}
Let $n$ and $k$ be natural numbers with $n>0$. The coefficient accuracy parameter in the
exact sparse-spectrum Goldreich–Levin input has value strictly less than half the
Fourier granularity at degree $k+1$.
\end{theorem}

\begin{theorem}[Fourier norm bounded by spectral sparsity]
\lean{FABL.fourierOneNorm_toReal_le_of_spectralSparsity_le}
\label{FABL.fourierOneNorm_toReal_le_of_spectralSparsity_le}
\leanok
\uses{FABL.BooleanFunction, FABL.BooleanFunction.toReal, FABL.binaryCubeSignEquiv, FABL.fourierOneNorm, FABL.fourierOneNorm_eq_spectralPNorm_one, FABL.signValue, FABL.spectralPNorm_one_eq_sum_abs, FABL.spectralSparsity, FABL.spectralSparsity_eq_card_vectorFourierSupport, FABL.sum_sq_vectorFourierCoeff_signValue_eq_one, FABL.vectorFourierCoeff, FABL.vectorFourierSupport}
\difficulty{5}
Let $n$ be a nonnegative integer, let $h$ be a sign-valued function on $\\{-1,1\\}^n$,
and let $S$ be a nonnegative integer. Regard $h$ as a real-valued function via its sign
values and as a function on the binary cube $\\mathbb F_2^n$ via the coordinatewise
identification of binary values with signs. If at most $S$ of its binary Fourier
coefficients are nonzero, then its Fourier one-norm on the sign cube is at most $S$.
\end{theorem}

\begin{theorem}[Granularity from Sparse Spectrum]
\lean{FABL.isFourierGranular_toReal_of_exactSparseSpectrumSparsity}
\label{FABL.isFourierGranular_toReal_of_exactSparseSpectrumSparsity}
\leanok
\uses{FABL.BooleanFunction, FABL.BooleanFunction.toReal, FABL.IsFourierGranular, FABL.IsVectorFourierGranular, FABL.binaryCubeSignEquiv, FABL.binaryFunctionOnSignCube, FABL.degreeFourierGranularity, FABL.exactSparseSpectrumSparsityBound, FABL.exactSparseSpectrumSparsityBound_two_le, FABL.isVectorFourierGranular_signValue_of_spectralSparsity_le, FABL.signValue, FABL.spectralSparsity, FABL.spectralSparsityGranularity_exactSparseSpectrumSparsityBound}
\difficulty{4}
Let $f:\{\pm1\}^n\to\{\pm1\}$, and let $k$ be a nonnegative integer. If the function on
$\mathbb F_2^n$ obtained by encoding the sign cube coordinates as bits and evaluating
$f$ has at most $2^{k+1}$ nonzero Walsh Fourier coefficients, then every Fourier
coefficient of $f$ is an integer multiple of $2^{-k}$.
\end{theorem}

\begin{theorem}[Granularity lower bound for nonzero coefficients]
\lean{FABL.degreeFourierGranularity_le_abs_fourierCoeff_of_exactSparseSpectrum}
\label{FABL.degreeFourierGranularity_le_abs_fourierCoeff_of_exactSparseSpectrum}
\leanok
\uses{FABL.BooleanFunction, FABL.BooleanFunction.toReal, FABL.binaryCubeSignEquiv, FABL.degreeFourierGranularity, FABL.degreeFourierGranularity_pos, FABL.exactSparseSpectrumSparsityBound, FABL.fourierCoeff, FABL.isFourierGranular_toReal_of_exactSparseSpectrumSparsity, FABL.signValue, FABL.spectralSparsity}
\difficulty{4}
Let $f$ be a Boolean function on the $n$-dimensional sign cube, and let $k$ be a natural
number. Suppose the number of nonzero Fourier coefficients of the corresponding function
on the binary cube is at most $2^{k+1}$. For any set $U$ of coordinates whose Fourier
coefficient for $f$ is nonzero, the absolute value of that coefficient is at least the
Fourier granularity at level $k+1$.
\end{theorem}

\begin{definition}[Rounded sparse Fourier hypothesis]
\lean{FABL.roundSparseSpectrumHypothesis}
\label{FABL.roundSparseSpectrumHypothesis}
\leanok
\uses{FABL.SparseFourierHypothesis, FABL.SparseFourierHypothesis.ofCoefficients, FABL.roundToDegreeFourierGranularity}
Given a natural number $k$ and a sparse Fourier hypothesis on dimension $n$, the rounded
hypothesis has the same character support, with each rational coefficient replaced by
its rounding to the Fourier granularity determined by $k+1$.
\end{definition}

\begin{theorem}[Support preserved by rounding]
\lean{FABL.roundSparseSpectrumHypothesis_support}
\label{FABL.roundSparseSpectrumHypothesis_support}
\leanok
\uses{FABL.SparseFourierHypothesis, FABL.roundSparseSpectrumHypothesis}
\difficulty{1}
For every natural number $k$ and every sparse Fourier hypothesis on $n$ coordinates, the
support of the rounded hypothesis equals the support of the original hypothesis.
\end{theorem}

\begin{theorem}[Rounded Fourier coefficient]
\lean{FABL.roundSparseSpectrumHypothesis_coefficient}
\label{FABL.roundSparseSpectrumHypothesis_coefficient}
\leanok
\uses{FABL.SparseFourierHypothesis, FABL.roundSparseSpectrumHypothesis, FABL.roundToDegreeFourierGranularity}
\difficulty{1}
For every natural number $k$, sparse Fourier hypothesis on $n$ coordinates, and
character $U$ in its support, the coefficient of $U$ after rounding the hypothesis is
the original coefficient rounded to the degree Fourier granularity at level $k+1$.
\end{theorem}

\begin{theorem}[Rounded Fourier coefficient identity]
\lean{FABL.roundedDeterministicGoldreichLevinHypothesis_coefficient_cast_eq}
\label{FABL.roundedDeterministicGoldreichLevinHypothesis_coefficient_cast_eq}
\leanok
\uses{FABL.BooleanFunction, FABL.BooleanFunction.toReal, FABL.DeterministicGoldreichLevinInput.coefficientInput, FABL.SmallBiasFourierInput.sample, FABL.abs_deterministicSmallBiasFourierEstimate_sub_fourierCoeff_le, FABL.binaryCubeSignEquiv, FABL.degreeFourierGranularity, FABL.deterministicGoldreichLevinHypothesisFromAnswers, FABL.exactSparseSpectrumGoldreichLevinInput, FABL.exactSparseSpectrumGoldreichLevinInput_coefficientEpsilon_lt_half_granularity, FABL.exactSparseSpectrumSparsityBound, FABL.fourierCoeff, FABL.fourierOneNorm, FABL.fourierOneNorm_toReal_le_of_spectralSparsity_le, FABL.isFourierGranular_toReal_of_exactSparseSpectrumSparsity, FABL.rationalSmallBiasFourierEstimate, FABL.rationalSmallBiasFourierEstimate_cast, FABL.roundSparseSpectrumHypothesis, FABL.roundToDegreeFourierGranularity, FABL.roundToDegreeFourierGranularity_eq_of_close, FABL.signValue, FABL.smallBiasFourierEstimate, FABL.spectralSparsity}
\difficulty{6}
Let $n$ be a positive integer, $k\ge 0$, and $f:\{\pm1\}^n\to\{\pm1\}$. Suppose the
function on $\mathbb F_2^n$ given by $x\mapsto f(x)$ under the coordinatewise
identification of the binary cube with the sign cube has at most $2^{k+1}$ nonzero
Fourier coefficients. For any finite family $\mathcal F$ of subsets of $\{1,\ldots,n\}$
and any $U\in\mathcal F$, form the empirical Fourier hypothesis on $\mathcal F$ from the
values of $f$ on the specified deterministic small-bias samples, then round each
coefficient to the nearest multiple of $2(1/2)^{k+1}$. The rounded coefficient at $U$
equals the real Fourier coefficient of $f$ indexed by $U$.
\end{theorem}

\begin{theorem}[Rounded hypothesis recovers the real target]
\lean{FABL.roundedDeterministicGoldreichLevinHypothesis_realValue_eq}
\label{FABL.roundedDeterministicGoldreichLevinHypothesis_realValue_eq}
\leanok
\uses{FABL.BooleanFunction, FABL.BooleanFunction.toReal, FABL.DeterministicGoldreichLevinInput.coefficientInput, FABL.DeterministicGoldreichLevinInput.threshold, FABL.SmallBiasFourierInput.sample, FABL.SparseFourierHypothesis.realValue, FABL.binaryCubeSignEquiv, FABL.degreeFourierGranularity, FABL.degreeFourierGranularity_le_abs_fourierCoeff_of_exactSparseSpectrum, FABL.deterministicGoldreichLevinHypothesisFromAnswers, FABL.exactSparseSpectrumGoldreichLevinInput, FABL.exactSparseSpectrumGoldreichLevinInput_threshold_le_granularity, FABL.exactSparseSpectrumSparsityBound, FABL.fourierCoeff, FABL.fourier_expansion, FABL.monomial, FABL.roundSparseSpectrumHypothesis, FABL.roundedDeterministicGoldreichLevinHypothesis_coefficient_cast_eq, FABL.signValue, FABL.spectralSparsity}
\difficulty{5}
Let $n$ be a positive integer, $k\ge 0$, and $f:\{\pm1\}^n\to\{\pm1\}$. Suppose the
function on $\mathbb F_2^n$ obtained through the coordinatewise binary-to-sign
identification has at most $2^{k+1}$ nonzero Fourier coefficients. Let $\mathcal F$ be a
finite family of coordinate subsets containing every subset $U$ whose real Fourier
coefficient for $f$ has magnitude at least the threshold determined by the exact
sparse-spectrum Goldreich–Levin input for $n,k$. Form the empirical Fourier hypothesis
from the target's values on the corresponding deterministic samples and round its
coefficients to granularity $2(1/2)^{k+1}$. The resulting real-valued Fourier sum equals
the real-valued function associated to $f$ at every point of the sign cube.
\end{theorem}

\begin{theorem}[Exact recovery from rounded Fourier estimates]
\lean{FABL.roundedDeterministicGoldreichLevinHypothesis_evaluate_eq}
\label{FABL.roundedDeterministicGoldreichLevinHypothesis_evaluate_eq}
\leanok
\uses{FABL.BooleanFunction, FABL.BooleanFunction.toReal, FABL.DeterministicGoldreichLevinInput.coefficientInput, FABL.DeterministicGoldreichLevinInput.threshold, FABL.SmallBiasFourierInput.sample, FABL.SparseFourierHypothesis.evaluate, FABL.SparseFourierHypothesis.evaluate_eq_thresholdSign_realValue, FABL.binaryCubeSignEquiv, FABL.deterministicGoldreichLevinHypothesisFromAnswers, FABL.exactSparseSpectrumGoldreichLevinInput, FABL.exactSparseSpectrumSparsityBound, FABL.fourierCoeff, FABL.roundSparseSpectrumHypothesis, FABL.roundedDeterministicGoldreichLevinHypothesis_realValue_eq, FABL.signValue, FABL.signValue_eq_neg_one_or_one, FABL.signValue_injective, FABL.spectralSparsity}
\difficulty{4}
Let $n$ be a positive integer, $k\in\mathbb N$, and $f:\{\pm1\}^n\to\{\pm1\}$. Suppose
the function on $\mathbb F_2^n$ obtained by identifying each binary coordinate with its
corresponding sign has at most $2^{k+1}$ nonzero Fourier coefficients. Let $\mathcal F$
be a family of subsets of $\{1,\ldots,n\}$ containing every $U$ whose Fourier
coefficient for $f$ has magnitude at least the threshold of the exact sparse-spectrum
Goldreich–Levin input with parameters $n,k$. Form the sparse Fourier hypothesis on
$\mathcal F$ from the empirical Fourier estimates of $f$ on the deterministic samples
specified by that input, and round each coefficient to granularity $2(1/2)^{k+1}$. The
sign-valued function obtained by assigning $+1$ to nonnegative hypothesis values and
$-1$ to negative values is equal to $f$ at every point of $\{\pm1\}^n$.
\end{theorem}

\begin{theorem}[Mapping preserves program cost]
\lean{FABL.DeterministicQueryProgram.runWithCost_map}
\label{FABL.DeterministicQueryProgram.runWithCost_map}
\leanok
\uses{FABL.DeterministicQueryProgram, FABL.DeterministicQueryProgram.map, FABL.DeterministicQueryProgram.runWithCost, FABL.DeterministicQueryProgram.runWithCost_bind}
\difficulty{3}
Let a deterministic query program have outputs in a type $A$, let $f:A\to B$ be any
function, and let an oracle assign an answer to each query. Running the program obtained
by applying $f$ to its output returns $f$ of the original run's output and exactly the
same learning cost.
\end{theorem}

\begin{definition}[Exact sparse-spectrum rounding work]
\lean{FABL.exactSparseSpectrumRoundingWork}
\label{FABL.exactSparseSpectrumRoundingWork}
\leanok
\uses{FABL.DeterministicGoldreichLevinInput, FABL.DeterministicGoldreichLevinInput.threshold, FABL.goldreichLevinActiveCap}
For a deterministic Goldreich–Levin input, the exact sparse-spectrum rounding work is
the active-cap value associated with that input's threshold.
\end{definition}

\begin{definition}[Positive exact sparse spectrum learner]
\lean{FABL.positiveExactSparseSpectrumLearner}
\statementsource{Analysis-of-Boolean-Functions-by-Ryan-ODonnell}{
  pdf-ffbe9c7b3306-p079-b001,
  pdf-ffbe9c7b3306-p083-b002
}
\label{FABL.positiveExactSparseSpectrumLearner}
\leanok
\uses{FABL.DeterministicQueryProgram, FABL.DeterministicQueryProgram.map, FABL.Sign, FABL.SparseFourierHypothesis, FABL.deterministicGoldreichLevinLearner, FABL.exactSparseSpectrumGoldreichLevinInput, FABL.exactSparseSpectrumRoundingWork, FABL.roundSparseSpectrumHypothesis}
For natural numbers $n$ and $k$ with $n>0$, consider the deterministic query program on
the $n$-dimensional Boolean cube that charges the exact sparse spectrum rounding work,
then runs the deterministic Goldreich–Levin learner with the corresponding exact sparse
spectrum input and rounds its output to a sparse Fourier hypothesis in dimension $n$.
\end{definition}

\begin{theorem}[Exact recovery by the sparse spectrum learner]
\lean{FABL.positiveExactSparseSpectrumLearner_evaluate_eq}
\label{FABL.positiveExactSparseSpectrumLearner_evaluate_eq}
\leanok
\uses{FABL.BooleanFunction, FABL.BooleanFunction.toReal, FABL.DeterministicGoldreichLevinInput.coefficientInput, FABL.DeterministicQueryProgram.map, FABL.DeterministicQueryProgram.runWithCost, FABL.DeterministicQueryProgram.runWithCost_bind, FABL.DeterministicQueryProgram.runWithCost_deterministicGoldreichLevinCoefficientProgram, FABL.DeterministicQueryProgram.runWithCost_map, FABL.SmallBiasFourierInput.sample, FABL.SparseFourierHypothesis.evaluate, FABL.binaryCubeSignEquiv, FABL.deterministicGoldreichLevinHypothesisFromAnswers, FABL.deterministicGoldreichLevinLearner, FABL.deterministicGoldreichLevinProgram, FABL.deterministicGoldreichLevinProgram_isCorrectOutput, FABL.deterministicGoldreichLevinSecondStage, FABL.exactSparseSpectrumGoldreichLevinInput, FABL.exactSparseSpectrumSparsityBound, FABL.fourierOneNorm, FABL.fourierOneNorm_toReal_le_of_spectralSparsity_le, FABL.positiveExactSparseSpectrumLearner, FABL.roundSparseSpectrumHypothesis, FABL.roundedDeterministicGoldreichLevinHypothesis_evaluate_eq, FABL.signValue, FABL.spectralSparsity}
\difficulty{5}
Let $n,k$ be natural numbers with $n>0$, and let $f:\\{-1,1\\}^n\\to\\{-1,1\\}$. Regard
$f$ as a function on $\\mathbb F_2^n$ using the coordinatewise binary-to-sign encoding.
If at most $2^{k+1}$ of its Fourier coefficients are nonzero, then running the positive
exact sparse spectrum learner with parameters $n,k$ and oracle $f$ produces a hypothesis
whose evaluation equals $f$ at every input.
\end{theorem}

\begin{definition}[Constant sparse Fourier hypothesis]
\lean{FABL.constantSparseFourierHypothesis}
\label{FABL.constantSparseFourierHypothesis}
\leanok
\uses{FABL.Sign, FABL.SparseFourierHypothesis, FABL.SparseFourierHypothesis.ofCoefficients}
For a dimension $n$ and a sign value, the constant sparse Fourier hypothesis has support
consisting only of the empty character, whose rational coefficient is the integer
represented by that sign.
\end{definition}

\begin{theorem}[Evaluation of a constant sparse Fourier hypothesis]
\lean{FABL.constantSparseFourierHypothesis_evaluate}
\label{FABL.constantSparseFourierHypothesis_evaluate}
\leanok
\uses{FABL.Sign, FABL.SparseFourierHypothesis.evaluate, FABL.SparseFourierHypothesis.evaluate_eq_thresholdSign_realValue, FABL.constantSparseFourierHypothesis, FABL.monomial, FABL.thresholdSign}
\difficulty{4}
For any natural number $n$ and any sign $b\in\{-1,1\}$, the Boolean function obtained by
evaluating the sparse Fourier hypothesis with its sole Fourier coefficient equal to $b$
is the constant function with value $b$.
\end{theorem}

\begin{definition}[Learner for zero-dimensional exact spectrum]
\lean{FABL.zeroDimensionExactSparseSpectrumLearner}
\label{FABL.zeroDimensionExactSparseSpectrumLearner}
\leanok
\uses{FABL.DeterministicQueryProgram, FABL.Sign, FABL.SparseFourierHypothesis, FABL.allOneSignCube, FABL.constantSparseFourierHypothesis}
For each dimension $n$, the learner makes one membership query at the all-one point of
the $n$-dimensional Boolean cube, then returns the constant sparse Fourier hypothesis
whose coefficient is the queried sign.
\end{definition}

\begin{theorem}[Cost of the zero-dimensional sparse-spectrum learner]
\lean{FABL.runWithCost_zeroDimensionExactSparseSpectrumLearner}
\label{FABL.runWithCost_zeroDimensionExactSparseSpectrumLearner}
\leanok
\uses{FABL.BooleanFunction, FABL.DeterministicQueryProgram.queryBatchCost, FABL.DeterministicQueryProgram.runWithCost, FABL.allOneSignCube, FABL.constantSparseFourierHypothesis, FABL.zeroDimensionExactSparseSpectrumLearner}
\difficulty{1}
For any natural number $n$ and Boolean target on the $n$-dimensional sign cube, the
learner that queries the target once at the all-one input and returns the constant
sparse Fourier hypothesis with that value has output equal to that hypothesis and cost
$(0,1,1)$.
\end{theorem}

\begin{theorem}[Evaluation of the zero-dimensional learner]
\lean{FABL.zeroDimensionExactSparseSpectrumLearner_evaluate_eq}
\label{FABL.zeroDimensionExactSparseSpectrumLearner_evaluate_eq}
\leanok
\uses{FABL.BooleanFunction, FABL.DeterministicQueryProgram.runWithCost, FABL.SparseFourierHypothesis.evaluate, FABL.constantSparseFourierHypothesis_evaluate, FABL.runWithCost_zeroDimensionExactSparseSpectrumLearner, FABL.zeroDimensionExactSparseSpectrumLearner}
\difficulty{4}
Let $n$ be a natural number with $\neg(0<n)$, and let $f$ be any Boolean function on the
$n$-dimensional sign cube. The hypothesis returned by running the zero-dimensional exact
sparse-spectrum learner with oracle $f$, when evaluated on the cube, agrees with $f$ at
every input.
\end{theorem}

\begin{definition}[Exact sparse spectrum learner]
\lean{FABL.exactSparseSpectrumLearner}
\statementsource{Analysis-of-Boolean-Functions-by-Ryan-ODonnell}{
  pdf-ffbe9c7b3306-p079-b001,
  pdf-ffbe9c7b3306-p083-b002,
  pdf-ffbe9c7b3306-p083-b003,
  pdf-ffbe9c7b3306-p159-b002
}
\label{FABL.exactSparseSpectrumLearner}
\leanok
\uses{FABL.DeterministicQueryProgram, FABL.Sign, FABL.SparseFourierHypothesis, FABL.positiveExactSparseSpectrumLearner, FABL.zeroDimensionExactSparseSpectrumLearner}
For natural numbers $n$ and $k$, use the positive exact sparse spectrum learner when
$n>0$, and use the zero-dimensional exact sparse spectrum learner when $n=0$. The
resulting program is a deterministic query program on the $n$-dimensional Boolean cube
with output a sparse Fourier hypothesis in dimension $n$.
\end{definition}

\begin{theorem}[Exact learner recovers the target]
\lean{FABL.exactSparseSpectrumLearner_evaluate_eq}
\label{FABL.exactSparseSpectrumLearner_evaluate_eq}
\leanok
\uses{FABL.BooleanFunction, FABL.DeterministicQueryProgram.runWithCost, FABL.SparseFourierHypothesis.evaluate, FABL.binaryCubeSignEquiv, FABL.exactSparseSpectrumLearner, FABL.exactSparseSpectrumSparsityBound, FABL.positiveExactSparseSpectrumLearner_evaluate_eq, FABL.signValue, FABL.spectralSparsity, FABL.zeroDimensionExactSparseSpectrumLearner_evaluate_eq}
\difficulty{4}
Let $n$ be a natural number, let $f:\{−1,1\}^n\to\{−1,1\}$ be a Boolean function, and
let $k$ be a natural number. If the number of nonzero Fourier coefficients of the
corresponding function on $\mathbb F_2^n$ is at most $2^{k+1}$, then the hypothesis
returned by the exact sparse spectrum learner, when evaluated on the sign cube, equals
$f$ at every input.
\end{theorem}

\begin{definition}[Exact sparse spectrum learning program]
\lean{FABL.exactSparseSpectrumLearningProgram}
\statementsource{Analysis-of-Boolean-Functions-by-Ryan-ODonnell}{
  pdf-ffbe9c7b3306-p079-b001,
  pdf-ffbe9c7b3306-p083-b002,
  pdf-ffbe9c7b3306-p159-b002
}
\label{FABL.exactSparseSpectrumLearningProgram}
\leanok
\uses{FABL.DeterministicQueryProgram.toLearningProgram, FABL.LearningProgram, FABL.SparseFourierHypothesis, FABL.exactSparseSpectrumLearner}
For natural numbers $n$ and $k$, the exact sparse spectrum learner with membership-query
access in dimension $n$ is a learning program whose output is a sparse Fourier
hypothesis on the $n$-dimensional Boolean cube.
\end{definition}

\begin{theorem}[Law of the exact sparse-spectrum learner]
\lean{FABL.runWithCost_exactSparseSpectrumLearningProgram}
\label{FABL.runWithCost_exactSparseSpectrumLearningProgram}
\leanok
\uses{FABL.BooleanFunction, FABL.DeterministicQueryProgram.runWithCost, FABL.DeterministicQueryProgram.runWithCost_toLearningProgram, FABL.LearningProgram.runWithCost, FABL.exactSparseSpectrumLearner, FABL.exactSparseSpectrumLearningProgram}
\difficulty{1}
Let $n,k\in\mathbb N$, and let $f:\{\pm1\}^n\to\{\pm1\}$ be a target function. The
outcome law of the exact sparse-spectrum learner is the point mass at the output and
learning cost obtained by running its deterministic query procedure against $f$.
\end{theorem}

\begin{definition}[Exact sparse spectrum query budget]
\lean{FABL.exactSparseSpectrumQueryBudget}
\label{FABL.exactSparseSpectrumQueryBudget}
\leanok
\uses{FABL.deterministicGoldreichLevinQueryBudget, FABL.exactSparseSpectrumGoldreichLevinInput}
For natural numbers $n$ and $k$, define the exact sparse spectrum query budget to be $1$
if $n=0$, and otherwise to be the deterministic Goldreich–Levin query budget of the
exact sparse-spectrum input with parameters $n$ and $k$.
\end{definition}

\begin{definition}[Exact sparse-spectrum work budget]
\lean{FABL.exactSparseSpectrumWorkBudget}
\label{FABL.exactSparseSpectrumWorkBudget}
\leanok
\uses{FABL.deterministicGoldreichLevinWorkBudget, FABL.exactSparseSpectrumGoldreichLevinInput, FABL.exactSparseSpectrumRoundingWork}
For natural numbers $n$ and $k$, define $B(n,k)$ as follows: if $n>0$, let $I$ be the
exact sparse-spectrum Goldreich–Levin input with parameters $n$ and $k$, and set
$B(n,k)$ equal to the sum of the rounding work for $I$ and the deterministic
Goldreich–Levin work budget for $I$; if $n=0$, set $B(n,k)=1$.
\end{definition}

\begin{theorem}[Exact learner resource bounds]
\lean{FABL.exactSparseSpectrumLearner_resource_bounds}
\label{FABL.exactSparseSpectrumLearner_resource_bounds}
\leanok
\uses{FABL.BooleanFunction, FABL.DeterministicQueryProgram.queryBatchCost, FABL.DeterministicQueryProgram.runWithCost, FABL.DeterministicQueryProgram.runWithCost_map, FABL.deterministicGoldreichLevinLearner, FABL.deterministicGoldreichLevinLearner_resource_bounds, FABL.deterministicGoldreichLevinQueryBudget, FABL.deterministicGoldreichLevinWorkBudget, FABL.exactSparseSpectrumGoldreichLevinInput, FABL.exactSparseSpectrumLearner, FABL.exactSparseSpectrumQueryBudget, FABL.exactSparseSpectrumRoundingWork, FABL.exactSparseSpectrumWorkBudget, FABL.positiveExactSparseSpectrumLearner, FABL.runWithCost_zeroDimensionExactSparseSpectrumLearner}
\difficulty{5}
For every natural number $n$, every Boolean function $f:\{−1,1\}^n\to\{−1,1\}$, and
every natural number $k$, running the exact sparse spectrum learner with oracle $f$ uses
zero random examples, makes at most the exact sparse spectrum query budget for $(n,k)$
membership queries, and performs at most the exact sparse spectrum work budget for
$(n,k)$ local computation steps.
\end{theorem}

\begin{definition}[Exact sparse spectrum runtime scale]
\lean{FABL.exactSparseSpectrumRuntimeScale}
\label{FABL.exactSparseSpectrumRuntimeScale}
\leanok
For natural numbers $n$ and $k$, define the exact sparse spectrum runtime scale by
$(n+1)(2^k+1)^2$.
\end{definition}

\begin{theorem}[Runtime scale bound]
\lean{FABL.exactSparseSpectrumGoldreichLevinInput_runtimeScale_le}
\label{FABL.exactSparseSpectrumGoldreichLevinInput_runtimeScale_le}
\leanok
\uses{FABL.DeterministicGoldreichLevinInput.learningParameter, FABL.DeterministicGoldreichLevinInput.learningParameter_cast, FABL.DeterministicGoldreichLevinInput.runtimeScale, FABL.SmallBiasInput.epsilon, FABL.exactSparseSpectrumGoldreichLevinInput, FABL.exactSparseSpectrumRuntimeScale, FABL.exactSparseSpectrumSparsityBound}
\difficulty{5}
Let $n$ and $k$ be natural numbers with $n>0$, and let $I$ be the deterministic
Goldreich–Levin input of dimension $n$, learning parameter $1/(4\cdot 2^k)$, and Fourier
bound $2^{k+1}$. Then the runtime scale of $I$ is at most $12(n+1)(2^k+1)^2$.
\end{theorem}

\begin{theorem}[Rounding work runtime bound]
\lean{FABL.exactSparseSpectrumRoundingWork_cast_le_goldreichLevinPolynomialRuntimeBound}
\label{FABL.exactSparseSpectrumRoundingWork_cast_le_goldreichLevinPolynomialRuntimeBound}
\leanok
\uses{FABL.DeterministicGoldreichLevinInput, FABL.DeterministicGoldreichLevinInput.activeCap_add_one_cast_le_runtimeScale, FABL.DeterministicGoldreichLevinInput.polynomialRuntimeBound, FABL.DeterministicGoldreichLevinInput.runtimeScale, FABL.DeterministicGoldreichLevinInput.runtimeScale_factor_bounds, FABL.DeterministicGoldreichLevinInput.threshold, FABL.exactSparseSpectrumRoundingWork, FABL.goldreichLevinActiveCap}
\difficulty{4}
Let $n$ be a positive integer, let $a,b$ be positive integers with $2a\le b$, and let
$s\ge1$ be an integer Fourier $1$-norm bound. Put $\varepsilon=a/b$ and
$R=(n+1)(s+1)/\varepsilon$. The exact sparse spectrum rounding work for these parameters
is at most $2^{170}R^{100}$.
\end{theorem}

\begin{definition}[Exact sparse spectrum polynomial runtime bound]
\lean{FABL.exactSparseSpectrumPolynomialRuntimeBound}
\statementsource{Analysis-of-Boolean-Functions-by-Ryan-ODonnell}{pdf-ffbe9c7b3306-p083-b002}
\label{FABL.exactSparseSpectrumPolynomialRuntimeBound}
\leanok
\uses{FABL.exactSparseSpectrumRuntimeScale}
For natural numbers $n$ and $k$, let $R(n,k)$ denote the exact sparse spectrum runtime
scale. The polynomial runtime bound is $1+2^{171}(12R(n,k))^{100}$.
\end{definition}

\begin{theorem}[Query budget runtime bound]
\lean{FABL.exactSparseSpectrumQueryBudget_cast_le_polynomialRuntimeBound}
\label{FABL.exactSparseSpectrumQueryBudget_cast_le_polynomialRuntimeBound}
\leanok
\uses{FABL.DeterministicGoldreichLevinInput.polynomialRuntimeBound, FABL.DeterministicGoldreichLevinInput.queryBudget_cast_le_polynomialRuntimeBound, FABL.DeterministicGoldreichLevinInput.runtimeScale, FABL.DeterministicGoldreichLevinInput.runtimeScale_factor_bounds, FABL.deterministicGoldreichLevinQueryBudget, FABL.exactSparseSpectrumGoldreichLevinInput, FABL.exactSparseSpectrumGoldreichLevinInput_runtimeScale_le, FABL.exactSparseSpectrumPolynomialRuntimeBound, FABL.exactSparseSpectrumQueryBudget, FABL.exactSparseSpectrumRuntimeScale}
\difficulty{5}
For all $n,k\in\mathbb{N}$, the exact sparse spectrum query budget, viewed as a real
number, is at most $1+2^{171}\bigl(12(n+1)(2^k+1)^2\bigr)^{100}$.
\end{theorem}

\begin{theorem}[Work budget runtime bound]
\lean{FABL.exactSparseSpectrumWorkBudget_cast_le_polynomialRuntimeBound}
\label{FABL.exactSparseSpectrumWorkBudget_cast_le_polynomialRuntimeBound}
\leanok
\uses{FABL.DeterministicGoldreichLevinInput.polynomialRuntimeBound, FABL.DeterministicGoldreichLevinInput.runtimeScale, FABL.DeterministicGoldreichLevinInput.runtimeScale_factor_bounds, FABL.DeterministicGoldreichLevinInput.workBudget_cast_le_polynomialRuntimeBound, FABL.deterministicGoldreichLevinWorkBudget, FABL.exactSparseSpectrumGoldreichLevinInput, FABL.exactSparseSpectrumGoldreichLevinInput_runtimeScale_le, FABL.exactSparseSpectrumPolynomialRuntimeBound, FABL.exactSparseSpectrumRoundingWork, FABL.exactSparseSpectrumRoundingWork_cast_le_goldreichLevinPolynomialRuntimeBound, FABL.exactSparseSpectrumRuntimeScale, FABL.exactSparseSpectrumWorkBudget}
\difficulty{5}
For all $n,k\in\mathbb{N}$, the exact sparse spectrum work budget, viewed as a real
number, is at most $1+2^{171}\bigl(12(n+1)(2^k+1)^2\bigr)^{100}$.
\end{theorem}

\begin{theorem}[Exact sparse spectrum learner queries polynomial le]
\lean{FABL.exactSparseSpectrumLearner_queries_polynomial_le}
\statementsource{Analysis-of-Boolean-Functions-by-Ryan-ODonnell}{
  pdf-ffbe9c7b3306-p083-b002,
  pdf-ffbe9c7b3306-p159-b002
}
\label{FABL.exactSparseSpectrumLearner_queries_polynomial_le}
\leanok
\uses{FABL.BooleanFunction, FABL.DeterministicQueryProgram.runWithCost, FABL.exactSparseSpectrumLearner, FABL.exactSparseSpectrumLearner_resource_bounds, FABL.exactSparseSpectrumPolynomialRuntimeBound, FABL.exactSparseSpectrumQueryBudget, FABL.exactSparseSpectrumQueryBudget_cast_le_polynomialRuntimeBound}
\difficulty{3}
% dataset-pipeline: stub -- prose pending build_informal_statements.py
Statement pending regeneration.
\end{theorem}

\begin{theorem}[Polynomial work bound for exact sparse learning]
\lean{FABL.exactSparseSpectrumLearner_work_polynomial_le}
\label{FABL.exactSparseSpectrumLearner_work_polynomial_le}
\leanok
\uses{FABL.BooleanFunction, FABL.DeterministicQueryProgram.runWithCost, FABL.exactSparseSpectrumLearner, FABL.exactSparseSpectrumLearner_resource_bounds, FABL.exactSparseSpectrumPolynomialRuntimeBound, FABL.exactSparseSpectrumWorkBudget, FABL.exactSparseSpectrumWorkBudget_cast_le_polynomialRuntimeBound}
\difficulty{3}
For every Boolean target function on the $n$-dimensional sign cube and every natural
number $k$, the local work charged when the exact sparse-spectrum learner is run against
the target is at most $1+2^{171}\bigl(12(n+1)(2^k+1)^2\bigr)^{100}$.
\end{theorem}

\begin{definition}[Exact sparse spectrum concept class]
\lean{FABL.exactSparseSpectrumConceptClass}
\statementsource{Analysis-of-Boolean-Functions-by-Ryan-ODonnell}{
  pdf-ffbe9c7b3306-p079-b001,
  pdf-ffbe9c7b3306-p159-b002
}
\label{FABL.exactSparseSpectrumConceptClass}
\leanok
\uses{FABL.BooleanFunction, FABL.binaryCubeSignEquiv, FABL.exactSparseSpectrumSparsityBound, FABL.signValue, FABL.spectralSparsity}
For natural numbers $n$ and $k$, the exact sparse spectrum concept class consists of the
sign-valued Boolean functions on the $n$-dimensional cube for which the spectral
sparsity of the associated $71$-valued function on the binary cube is at most
$2^{k+1}$.
\end{definition}

\begin{definition}[Exact sparse-spectrum learning algorithm]
\lean{FABL.exactSparseSpectrumLearningAlgorithm}
\statementsource{Analysis-of-Boolean-Functions-by-Ryan-ODonnell}{
  pdf-ffbe9c7b3306-p079-b001,
  pdf-ffbe9c7b3306-p083-b002,
  pdf-ffbe9c7b3306-p159-b002
}
\label{FABL.exactSparseSpectrumLearningAlgorithm}
\leanok
\uses{FABL.LearningAlgorithm, FABL.LearningCost, FABL.LearningProgram.eventProbability, FABL.SparseFourierHypothesis, FABL.SparseFourierHypothesis.evaluate, FABL.SparseFourierHypothesis.finiteRepresentation, FABL.exactSparseSpectrumLearner_resource_bounds, FABL.exactSparseSpectrumLearningProgram, FABL.exactSparseSpectrumQueryBudget, FABL.exactSparseSpectrumWorkBudget, FABL.relativeHammingDist, FABL.runWithCost_exactSparseSpectrumLearningProgram}
For each dimension $n$ and sparsity parameter $k$, the construction assigns a
membership-query learning algorithm for sparse Fourier hypotheses. For every requested
accuracy $psilon\in[0,1/2]$, its program is the exact sparse-spectrum learner for $n$
and $k$; on every Boolean target, its advertised success probability is the probability
that the output hypothesis has relative Hamming distance at most $psilon$ from the
target. Every execution uses zero random examples, and its membership queries and local
work are bounded by the corresponding query and work budgets for that learner.
\end{definition}

\begin{theorem}[Exact sparse spectrum learning algorithm success probability eq one]
\lean{FABL.exactSparseSpectrumLearningAlgorithm_successProbability_eq_one}
\statementsource{Analysis-of-Boolean-Functions-by-Ryan-ODonnell}{pdf-ffbe9c7b3306-p159-b002}
\label{FABL.exactSparseSpectrumLearningAlgorithm_successProbability_eq_one}
\leanok
\uses{FABL.BooleanFunction, FABL.LearningAccuracy, FABL.LearningProgram.eventProbability, FABL.SparseFourierHypothesis.evaluate, FABL.exactSparseSpectrumConceptClass, FABL.exactSparseSpectrumLearner_evaluate_eq, FABL.exactSparseSpectrumLearningAlgorithm, FABL.exactSparseSpectrumLearningProgram, FABL.relativeHammingDist, FABL.runWithCost_exactSparseSpectrumLearningProgram}
\difficulty{3}
% dataset-pipeline: stub -- prose pending build_informal_statements.py
Statement pending regeneration.
\end{theorem}

\begin{theorem}[Exact sparse spectrum learning algorithm learns]
\lean{FABL.exactSparseSpectrumLearningAlgorithm_learns}
\statementsource{Analysis-of-Boolean-Functions-by-Ryan-ODonnell}{
  pdf-ffbe9c7b3306-p159-b002,
  pdf-ffbe9c7b3306-p083-b002,
  pdf-ffbe9c7b3306-p083-b003,
  pdf-ffbe9c7b3306-p079-b001
}
\label{FABL.exactSparseSpectrumLearningAlgorithm_learns}
\leanok
\uses{FABL.LearningAccuracy, FABL.LearnsConceptClassWithError, FABL.exactSparseSpectrumConceptClass, FABL.exactSparseSpectrumLearningAlgorithm, FABL.exactSparseSpectrumLearningAlgorithm_successProbability_eq_one}
\difficulty{2}
% dataset-pipeline: stub -- prose pending build_informal_statements.py
Statement pending regeneration.
\end{theorem}

\begin{theorem}[Exact learning for sparse spectra]
\lean{FABL.exactSparseSpectrumLearningProgram_zero_error}
\proofsource{Analysis-of-Boolean-Functions-by-Ryan-ODonnell}{pdf-ffbe9c7b3306-p159-b002}
\label{FABL.exactSparseSpectrumLearningProgram_zero_error}
\leanok
\uses{FABL.BooleanFunction, FABL.LearningCost, FABL.LearningProgram.runWithCost, FABL.SparseFourierHypothesis, FABL.SparseFourierHypothesis.evaluate, FABL.exactSparseSpectrumConceptClass, FABL.exactSparseSpectrumLearner_evaluate_eq, FABL.exactSparseSpectrumLearningProgram, FABL.relativeHammingDist, FABL.runWithCost_exactSparseSpectrumLearningProgram}
\difficulty{3}
Let $n,k\in\mathbb{N}$ and let $f:\{\pm1\}^n\to\{\pm1\}$ have at most $2^{k+1}$ nonzero
Fourier coefficients. Every hypothesis returned with positive probability by the exact
sparse spectrum learning program on target $f$ agrees with $f$ at every input; in
particular, its relative Hamming distance from $f$ is zero.
\end{theorem}

\begin{theorem}[Exact sparse-spectrum learning guarantee]
\lean{FABL.exactSparseSpectrumLearningProgram_spec}
\proofsource{Analysis-of-Boolean-Functions-by-Ryan-ODonnell}{pdf-ffbe9c7b3306-p159-b002}
\label{FABL.exactSparseSpectrumLearningProgram_spec}
\leanok
\uses{FABL.BooleanFunction, FABL.DeterministicQueryProgram.runWithCost, FABL.LearningCost, FABL.LearningProgram.runWithCost, FABL.SparseFourierHypothesis, FABL.SparseFourierHypothesis.evaluate, FABL.exactSparseSpectrumConceptClass, FABL.exactSparseSpectrumLearner, FABL.exactSparseSpectrumLearner_resource_bounds, FABL.exactSparseSpectrumLearningProgram, FABL.exactSparseSpectrumLearningProgram_zero_error, FABL.exactSparseSpectrumPolynomialRuntimeBound, FABL.exactSparseSpectrumQueryBudget, FABL.exactSparseSpectrumQueryBudget_cast_le_polynomialRuntimeBound, FABL.exactSparseSpectrumWorkBudget, FABL.exactSparseSpectrumWorkBudget_cast_le_polynomialRuntimeBound, FABL.relativeHammingDist, FABL.runWithCost_exactSparseSpectrumLearningProgram}
\difficulty{4}
Let $n$ and $k$ be natural numbers, and let the target be a Boolean function on the
$n$-dimensional sign cube whose Fourier spectrum has at most $2^{k+1}$ nonzero
coefficients. For every outcome in the support of the learning program’s cost
distribution on this target, the output hypothesis evaluates exactly to the target and
has relative Hamming distance zero from it; the outcome uses zero random examples, and
both its number of membership queries and its charged local work are at most
$1+2^{171}\bigl(12(n+1)(2^k+1)^2\bigr)^{100}$.
\end{theorem}
